\label{sec:alabama}

\subsection{Background: \textit{Allen v.\ Milligan}}

\textit{Allen v.\ Milligan}~\citep{allen2023} arose from Alabama's 2021
congressional redistricting, in which the legislature drew a map with only
one majority-Black congressional district (the 7th, in the Birmingham-to-Tuscaloosa
corridor) despite Black voters constituting 27\% of Alabama's voting-age
population.  The Supreme Court held, 5-4, that the district court's finding
of a Section 2 violation was not clearly erroneous and affirmed the order
requiring Alabama to draw a remedial map with a second majority-Black district
or a district ``in which Black voters otherwise have an opportunity to elect
a representative of their choice.''

We demonstrate how the algorithm's VRASection output would serve as the
remedial plan in the post-\textit{Allen v.\ Milligan} (2023) remedial
context: the court ordered Alabama to draw a second majority-Black
district, making Alabama the canonical post-\textit{Allen} test case.
Alabama provides an ideal worked example for our Gingles methodology because:
(a) the three prongs have been litigated exhaustively in the district court
record; (b) the Black community in the Black Belt region of Alabama is
geographically concentrated in a way that makes Prong 1 straightforward; and
(c) the post-\textit{Callais} Prong 3 disentanglement is particularly
important in Alabama, where Black voters are overwhelmingly Democratic, making
the racial/partisan distinction contested.

\subsection{Prong 1: Geographic Compactness}

Running \texttt{redist state --state AL --year 2020 --version remedy --vra-mode}
generates an algorithmic map with two majority-Black districts: the existing
7th district (Birmingham corridor) and a new district traversing the Black Belt
counties (Wilcox, Dallas, Lowndes, Marengo, Greene, Hale, Perry) from the
Georgia border toward Mobile.

The VRASection alignment scores for the two districts are:
\begin{itemize}
  \item \textbf{7th District (Birmingham corridor):} $A_d = 0.73$, Black CVAP
    = 55.2\%.  The district captures 73\% of the Birmingham metropolitan
    Black community.  Prong 1 is satisfied.
  \item \textbf{New Black Belt district:} $A_d = 0.61$, Black CVAP = 50.4\%.
    The district captures 61\% of the Black Belt minority community.  Prong 1
    is satisfied under the 0.5 threshold.
\end{itemize}

The geographic visualization generated by \texttt{redist map} shows that the
new Black Belt district follows county lines through the Black Belt region,
with the irregular geometry attributable to population-balance requirements
rather than racial gerrymandering.

\subsection{Prong 2: Minority Political Cohesion}

We analyze five elections in the Alabama 7th district area: the 2018 Democratic
Senate primary (Jones vs.\ competition), the 2020 Democratic Senate primary
(run-off), the 2018 congressional primary in the 7th, the 2020 congressional
primary in the 7th, and the 2022 congressional primary in the 7th.

The WLS regression produces $\hat{\beta}_{\text{min}}$ estimates ranging from
0.84 to 0.91 across the five elections.  Black voters in Alabama support
Black-preferred candidates at rates between 84\% and 91\%, a pattern
consistent with strong political cohesion.  The HC3 standard errors range
from 0.03 to 0.05, and all five estimates are statistically significant at
the $p < 0.001$ level.  Prong 2 is satisfied.

\subsection{Prong 3: Majority Bloc Voting with Partisan Control}

We analyze eight general elections in Alabama congressional districts: the
2014, 2016, 2018, 2020, and 2022 congressional elections in the 7th district
and adjacent districts, plus the 2017 and 2020 Senate general elections.

The baseline regression without partisan control produces $\hat{\beta}_{\text{race}}
= 0.71$ (SE = 0.04), indicating that majority voters support the candidate
opposed by Black voters at a 71\% rate.  However, when the partisan control
variable (Republican presidential vote share) is included, the racial component
drops to $\hat{\beta}_{\text{race}} = 0.43$ (SE = 0.06).  The remaining 0.43
represents the racial component of majority bloc voting that is not attributable
to partisan alignment.

\begin{table}[h]
\centering
\caption{Prong 3 WLS+HC3 Results, Alabama (Selected Elections)}
\label{tab:alabama-prong3}
\begin{tabular}{lcccc}
\toprule
Election & $\hat{\beta}_{\text{race}}$ (no control) & $\hat{\beta}_{\text{race}}$ (w/ control) & HC3 SE & Holm-$p$ \\
\midrule
2020 AL-07 General & 0.74 & 0.45 & 0.07 & $<$0.001 \\
2022 AL-07 General & 0.69 & 0.41 & 0.08 & $<$0.001 \\
2020 Senate General & 0.73 & 0.44 & 0.06 & $<$0.001 \\
2017 Senate General & 0.71 & 0.42 & 0.07 & $<$0.001 \\
2018 AL-02 General & 0.68 & 0.39 & 0.09 & 0.001 \\
\bottomrule
\end{tabular}
\end{table}

The Holm correction across all eight elections yields corrected $p$-values
below 0.001 for five elections and below 0.01 for all eight.  The LOO analysis
shows that excluding any single election does not change the Prong 3 conclusion.
Prong 3 is satisfied under the post-\textit{Callais} standard: majority
bloc voting has a racial component with a point estimate of 0.41--0.45,
independent of partisan alignment, and this estimate is stable, statistically
significant, and robust to the multiple-testing correction.

\subsection{Summary: All Three Prongs Satisfied for Alabama AL-07 and New Black Belt District}

The worked example demonstrates that both the existing 7th district and the
new Black Belt district satisfy all three Gingles prongs under the methodology
developed in Sections~2--4.  The quantitative outputs---alignment scores,
WLS cohesion estimates, racial bloc-voting coefficients with partisan
controls---provide the evidentiary basis for expert testimony without relying
on visual inspection or ad hoc judgment.
